\section*{Результаты работы}

\subsection*{Сборка и запуск}
\begin{verbatim}
PS C:\OSi\lab2> mingw32-make
gcc -Wall -Wextra -Iinc -o parent.exe src/parent.c src/common.c src/platform_win.c
gcc -Wall -Wextra -Iinc -o child.exe src/child.c src/common.c
PS C:\OSi\lab2> .\parent.exe
Enter filename: result.txt
Enter numbers (space-separated): 1 5 4 2 4 2.3
Parent received sum: 18.30
Child finished.
\end{verbatim}

\subsection*{Содержимое файла result.txt}
\begin{verbatim}
Sum: 18.30
\end{verbatim}

\subsection*{Диагностика через Process Monitor}
Для подтверждения работы системных вызовов использовался Process Monitor (ProcMon). В логах зафиксированы события:
\begin{itemize}
    \item \texttt{CreateFile} — создание анонимных каналов.
    \item \texttt{Process Create} — запуск дочернего процесса \texttt{child.exe}.
    \item \texttt{WriteFile} — запись чисел родителем в канал.
    \item \texttt{ReadFile} — чтение результата родителем из канала.
\end{itemize}